\documentclass[10pt,a4paper]{amsart}
\usepackage[margin=19mm]{geometry}
\usepackage{amsmath,amssymb,mathtools}
\usepackage[scr=rsfs,cal=euler]{mathalfa}
\usepackage{xfrac}
\usepackage[unicode,hidelinks,bookmarks=false]{hyperref}
\newcommand{\ie}{i.e.\ }
\newcommand{\cf}{cf.\ }
\newcommand{\resp}{respectively\ }
\newcommand{\Subsection}{\subsection}
\providecommand{\nobreakdash}{\nobreak\hbox{-}}
\setlength{\emergencystretch}{2em}
\multlinegap0pt
\usepackage{bm}
\usepackage{bbm}
\usepackage{dsfont}
\usepackage{xspace}
\usepackage{cancel}
\usepackage{mathtools}
\usepackage{amsthm}
\makeatletter
\@ifundefined{lemma}{\newtheorem{lemma}{Lemma}}{}
\renewenvironment{proof}[1]{\par\medskip\noindent\textbf{#1}\enspace\ignorespaces}{\par\medskip}
\makeatother
\def\E{\mathbb{E}}%
\def\P{\mathbb{P}}
\title[Wasserstein-p Central Limit Theorem Rates: From Local Depend]{Wasserstein-p Central Limit Theorem Rates: From Local Dependence to Markov Chains}
\author{Yixuan Zhang, Qiaomin Xie}
\date{Source: arXiv:2601.08184v3 (CC BY 4.0)}
\begin{document}
\maketitle

\noindent\emph{Source and licence.} Wasserstein-p Central Limit Theorem Rates: From Local Dependence to Markov Chains. Version arXiv:2601.08184v3 (CC BY 4.0), distributed under the Creative Commons Attribution 4.0 International licence, as recorded in the arXiv licence metadata for this exact version and shown on the abstract page https://arxiv.org/abs/2601.08184v3. Full rights notice: licenses/arxiv-2601.08184-CC-BY-4.0.txt.\par\smallskip

\noindent\textit{Editorial scope.} Counted unit: the Lemma whose source label is \texttt{lem:K} in the article (source0.tex lines 2900--2908, source label \texttt{lem:K}), delivered together with the article's own complete original proof of it (source0.tex lines 2909--3033, one \texttt{proof} environment of 125 lines). No mathematics is written, repaired or re-derived: statement and proof are the article's own text line for line. Required context retained uncounted: lem:geom\_tail (lines 1310--1316). Every definition, display and label of those lines is retained unchanged. Omitted from this package and not redistributed: 1--1309, 1317--2899, 3034--3973. Nothing inside a retained excerpt is omitted or altered: every retained body is delivered exactly as the article's lines are written, so no omission receipt of any kind accompanies this package. Citations: the retained excerpts carry no citation command at all, so no bibliography entry of the article is transcribed and no reference is redistributed. Reference adaptation: none was needed and none was made; every reference inside the delivered unit resolves inside it, so no unresolved reference or citation is produced. Printed slips kept literal: no printed word, symbol, hypothesis or number of the counted statement or of its proof was altered. Layout and class: the article's own class is \texttt{informs3} with packages including theorem, tgtermes, newtxtext, newtxmath, bm, endnotes, inputenc, enumitem, ulem, hyperref, booktabs, xcolor, todonotes, color-edits; the wrapper is the lane's pinned \texttt{amsart} editorial wrapper at 10pt on A4, which restates the theorem-like environments the retained text needs and adds the article's own macro definitions that the retained text needs (\texttt{\textbackslash E}, \texttt{\textbackslash P}), with extra wrapper packages (bm, bbm, dsfont, xspace, cancel, mathtools). The delivered numbering is the unit's own. Assets: no retained span contains a figure environment, a graphics-inclusion command, a PDF-page inclusion, a table or a TikZ picture; no image file is copied into this package, the package declares no asset dependency and no image file is redistributed with it. Licence: version arXiv:2601.08184v3 of this article is distributed under the Creative Commons Attribution 4.0 International licence, as recorded in the arXiv licence metadata for this exact version and shown on the abstract page https://arxiv.org/abs/2601.08184v3. A full rights notice is delivered at licenses/arxiv-2601.08184-CC-BY-4.0.txt. Characters: every retained excerpt is pure ASCII, so the delivered engine, XeLaTeX with its default Latin Modern text font, reports no missing character.
\begin{lemma}\label{lem:geom_tail}
The cycle lengths \(\{L_i\}_{i\ge 1}\) are independent, \(\{L_i\}_{i\ge 2}\) are i.i.d., and there exist constants \(b>0\) and \(\rho\in(0,1)\) such that
\[
\mathbb P(L_i>\ell)\le b\rho^\ell,
\qquad \forall i\ge 1,\ \ell\ge 0.
\]
\end{lemma}

\begin{lemma}\label{lem:K}
Under Lemma~\ref{lem:geom_tail}, for every \(q\ge 1\) there exists
\(C_q<\infty\) such that
\[
\E\bigl|K_n-\lfloor n/\E[\Lambda_2] \rfloor\bigr|^{q}
\le C_q\, n^{q/2},
\qquad \forall n\ge 1.
\]
\end{lemma}

\begin{proof}{Proof of Lemma \ref{lem:K}}Write
\[
\mu:=\E[\Lambda_2],\qquad
k_0:=\lfloor n/\mu\rfloor,
\qquad\text{and}\qquad
N_n:=\min\{k\ge1:T_k>n\},
\]
so that $K_n=N_n+1.$ We first show that the cycle lengths admit a uniform exponential moment.
By Lemma~\ref{lem:geom_tail}, there exist constants \(b<\infty\) and
\(\rho\in(0,1)\) such that
\[
\sup_{i\ge1}\P(L_i>\ell)\le b\,\rho^\ell,
\qquad \forall \ell\ge0.
\]
Fix \(\theta_L\in(0,-\log\rho)\), and set
\[
r:=e^{\theta_L}\rho\in(0,1).
\]
For any nonnegative integer-valued random variable \(Y\), we have the identity
\[
\E[e^{\theta_L Y}]
=
1+\sum_{\ell\ge0}\bigl(e^{\theta_L(\ell+1)}-e^{\theta_L\ell}\bigr)\P(Y>\ell).
\]
Applying this with \(Y=L_i\), we obtain
\[
\E[e^{\theta_L L_i}]
=
1+(e^{\theta_L}-1)\sum_{\ell\ge0}e^{\theta_L\ell}\P(L_i>\ell)
\le
1+(e^{\theta_L}-1)b\sum_{\ell\ge0}(e^{\theta_L}\rho)^\ell.
\]
Since \(r<1\), the geometric series converges, and therefore $\sup_{i\ge1}\E[e^{\theta_L L_i}]<\infty.$ Because \(\Lambda_i=mL_i\), it follows that there exists
\(\theta_\Lambda>0\) such that
\begin{equation}\label{eq:Lambda-exp-moment}
\sup_{i\ge1}\E[e^{\theta \Lambda_i}]<\infty,
\qquad \forall \theta\in(0,\theta_\Lambda].
\end{equation}
Next, by Lemma \ref{lem:geom_tail}, the variables \[ X_i:=\Lambda_i-\mu,\qquad i\ge 1, \] are independent and sub-exponential by \eqref{eq:Lambda-exp-moment} and centered when $i \geq 2$. Therefore, Bernstein's inequality yields constants \(c,C>0\) such that, for all \(k\ge 1\) and \(u\ge0\), \[ \P\Bigl(\Bigl|\sum_{i=1}^{k}X_i\Bigr|\ge u\Bigr) \le C\exp\!\left[-c\min\!\left(\frac{u^2}{k},u\right)\right]. \] Since $T_k-k\mu = \sum_{i=1}^{k}X_i,$ \begin{equation}\label{eq:Tk_conc} \P\bigl(|T_k-k\mu|\ge u\bigr) \le C\exp\!\left[-c\min\!\left(\frac{u^2}{k},u\right)\right]. \end{equation}
We now derive a deviation bound for \(N_n\). For any integer \(r\ge2\),
\[
\{N_n\ge k_0+r\}=\{T_{k_0+r-1}\le n\}.
\]
Moreover,
\[
n-(k_0+r-1)\mu
=
(n-k_0\mu)-(r-1)\mu
\le
\mu-(r-1)\mu
=
-(r-2)\mu,
\]
and hence
\[
\{N_n\ge k_0+r\}
\subseteq
\Bigl\{
T_{k_0+r-1}-(k_0+r-1)\mu\le -(r-2)\mu
\Bigr\}.
\]
Similarly, for \(1\le r\le k_0-1\),
\[
\{N_n\le k_0-r\}=\{T_{k_0-r}>n\},
\]
and since
\[
n-(k_0-r)\mu=(n-k_0\mu)+r\mu\ge r\mu,
\]
we obtain
\[
\{N_n\le k_0-r\}
\subseteq
\Bigl\{
T_{k_0-r}-(k_0-r)\mu\ge r\mu
\Bigr\}.
\]
When \(r\ge k_0\), the event \(\{N_n\le k_0-r\}\) is empty. Applying \eqref{eq:Tk_conc} and using that
\(k_0+r-1\lesssim n+r\) and \(k_0-r\lesssim n\), we obtain constants
\(c',C'>0\) such that, for all \(r\ge1\),
\begin{equation}\label{eq:N_tail}
\P\bigl(|N_n-k_0|\ge r\bigr)
\le
C'\exp\!\left[-c'\min\!\left(\frac{r^2}{n},r\right)\right].
\end{equation}
Finally, by the tail-integral formula,
\[
\E|N_n-k_0|^q
=
q\int_0^\infty t^{q-1}\P(|N_n-k_0|\ge t)\,dt.
\]
Using \eqref{eq:N_tail}, we get
\[
\E|N_n-k_0|^q
\le
qC'\int_0^\infty t^{q-1}
\exp\!\left[-c'\min\!\left(\frac{t^2}{n},t\right)\right]dt.
\]
We split the integral at \(t=n\). For \(0\le t\le n\), we have
\(\min(t^2/n,t)=t^2/n\), and the change of variables \(t=\sqrt n\,x\) yields
\[
\int_0^n t^{q-1}e^{-c't^2/n}\,dt
\lesssim
n^{q/2}\int_0^\infty x^{q-1}e^{-c'x^2}\,dx
\lesssim
n^{q/2}.
\]
For \(t\ge n\), we have \(\min(t^2/n,t)=t\), so $\int_n^\infty t^{q-1}e^{-c't}dt<\infty.$ Thus 
\[
\E|N_n-k_0|^q\le \widetilde C_q n^{q/2}
\]
for some constant \(\widetilde C_q<\infty\). Since \(K_n=N_n+1\), we have
\[
|K_n-k_0|^q\le 2^{q-1}\bigl(|N_n-k_0|^q+1\bigr),
\]
and therefore
\[
\E|K_n-k_0|^q
\le
2^{q-1}\bigl(\E|N_n-k_0|^q+1\bigr)
\le
C_q\, n^{q/2}.
\]
This completes the proof.
\end{proof}

\end{document}
